
\section{Token Requirements}\label{sec:token-requirements}

\textit{Normative language per \href{https://tools.ietf.org/html/rfc2119}{RFC~2119}.}

\subsection{Overview}


This section defines the complete structural requirements for JWT access tokens issued
by a tenant's IdP and presented to the Canvasflow GraphQL API. A token MUST satisfy
every requirement in \hyperref[sec:signing-path-requirements]{\S5.2}\ through 
\hyperref[sec:prohibited]{\S5.5} to be accepted. 
% The normative validation sequence the Resource Server applies is defined in \texttt{spec\_appendix\_validation}.

All example values in this section use: \linebreak
\texttt{iss = https://login.tenant.com}, 
\texttt{aud = canvasflow-api}, and entitlement values in the mandatory
\texttt{cf:<type>:<id>} format.

\subsection{Signing Path Requirements}\label{sec:signing-path-requirements}

Canvasflow supports two signing paths. The choice MUST be made at onboarding and
recorded in the tenant configuration. The two paths are not interchangeable after go-live
without a coordinated migration.

\subsubsection{Path A --- JWKS Endpoint (Preferred)}\label{sec:jwks-endpoint}

The tenant's IdP signs access tokens using an asymmetric algorithm (RS256 or ES256) and
publishes the corresponding public keys at a stable JSON Web Key Set (JWKS) endpoint per
\href{https://datatracker.ietf.org/doc/html/rfc7517}{RFC~7517}.

\textbf{Requirements:}

\begin{itemize}
  \item Algorithm MUST be RS256 or ES256 (RFC~7518~\href{https://datatracker.ietf.org/doc/html/rfc7518\#section-3.3}{\S3.3} and \href{https://datatracker.ietf.org/doc/html/rfc7518\#section-3.4}{\S3.4}).
  \item RSA keys MUST be at minimum 2048~bits in length.
  \item EC keys MUST use the P-256 curve.
  \item The \texttt{kid} header claim MUST be present in every issued token and MUST
    match a key in the JWKS endpoint.
  \item The JWKS URL MUST be publicly accessible (no authentication required to fetch it).
  \item The tenant MUST communicate the exact JWKS URL to Canvasflow at onboarding.
  \item \textbf{Key rotation:} the tenant MUST publish the new key alongside the outgoing
    key for a minimum overlap window of one hour before removing the outgoing key.
    Canvasflow resolves \texttt{kid} dynamically against a one-hour cache; an unknown
    \texttt{kid} triggers an immediate re-fetch. The one-hour overlap window is required
    to prevent transient 401 errors during rotation.
\end{itemize}

\begin{table}[htbp]
\centering
\begin{tabularx}{\textwidth}{@{}lX@{}}
\toprule
\textbf{IdP} & \textbf{JWKS URL format} \\
\midrule
AWS Cognito      & {\scriptsize\texttt{https://cognito-idp.<region>.amazonaws.com/<pool-id>/.well-known/jwks.json}} \\[4pt]
Okta             & {\scriptsize\texttt{https://<tenant>.okta.com/oauth2/v1/keys}} \\[4pt]
Azure AD / Entra & {\scriptsize\texttt{https://login.microsoftonline.com/<tenant-id>/discovery/v2.0/keys}} \\[4pt]
Auth0            & {\scriptsize\texttt{https://<tenant>.auth0.com/.well-known/jwks.json}} \\[4pt]
\bottomrule
\end{tabularx}
\caption{Well-known JWKS endpoint formats for major IdPs}
\end{table}

\pagebreak

\subsubsection{Path B --- HS256 Shared Secret (Exception)}

The tenant's IdP signs access tokens using HMAC-SHA256 with a shared secret that both
the tenant and Canvasflow hold. This path is an explicit exception reserved for custom
or lightweight IdPs that cannot publish a JWKS endpoint. It MUST be explicitly approved
and documented by Canvasflow before use.

\textbf{Requirements:}

\begin{itemize}
  \item Algorithm MUST be HS256.
  \item The shared secret MUST be at minimum 32~bytes (256~bits) of cryptographically
    random data. A secret shorter than 32~bytes MUST be rejected.
  \item The \texttt{kid} header claim MUST be present in every issued token and MUST
    match a known secret version in Canvasflow's configuration. This supports rotation.
  \item The tenant MUST share the secret with Canvasflow via a secure channel --- never
    via unencrypted email.
  \item \textbf{Key rotation:} the tenant MUST notify Canvasflow before rotating the
    secret, provide the new secret via a secure channel, and maintain the outgoing secret
    as valid under its current \texttt{kid} for at minimum one hour after the new secret
    begins signing tokens.
\end{itemize}

\textbf{Justification for JWKS over HS256.} The following four considerations make JWKS
the strongly preferred path:

\begin{enumerate}
  \item \textbf{Asymmetry of forgery risk.} With JWKS (RS256/ES256), Canvasflow holds
    only the tenant's public key. A total compromise of Canvasflow's infrastructure gives
    an attacker no ability to forge valid tokens for any JWKS tenant. With HS256, the
    shared secret is simultaneously a verification key and a signing key: a breach of
    Canvasflow's secret store allows an attacker to forge arbitrary valid tokens for
    every HS256 tenant.

  \item \textbf{Self-service key rotation.} A JWKS tenant rotates keys by publishing a
    new key pair under a new \texttt{kid} alongside the outgoing key. Canvasflow
    resolves \texttt{kid} dynamically; rotation requires zero coordination with Canvasflow
    and incurs zero downtime. HS256 rotation requires out-of-band secret exchange,
    coordinated timing, and a Canvasflow-side configuration update on every rotation.

  \item \textbf{No secret-handling liability.} A JWKS URL is public information. Every
    HS256 secret Canvasflow holds is standing security liability (\hyperref[sec:canvasflow-guarantees]{\S11}).

  \item \textbf{Zero marginal cost for enterprise IdPs.} AWS Cognito, Okta, Azure AD /
    Entra, and Auth0 all publish JWKS endpoints by default.
\end{enumerate}

\textbf{Conclusion:} JWKS is the default. HS256 is an explicit, documented exception
accepted as a trade of security posture for reduced onboarding friction, available only
when the tenant has confirmed that a JWKS endpoint is not achievable.

\subsection{Required Header Claims}

The following claims MUST be present in the JWT header of every access token:



\subsubsection{\texttt{typ} --- Token Type (RFC~9068~\S2.1)}\label{sec:token-type}

The \texttt{typ} header MUST be set to exactly \texttt{"at+jwt"} or
\texttt{"application/at+jwt"}.

\textbf{Rationale.} This is the primary defense against token confusion attacks. An
attacker who obtains an OIDC ID token --- which is signed by the same key as the access
token and carries valid \texttt{iss}, \texttt{sub}, and \texttt{aud} claims --- could
present it to the GraphQL API. Without \texttt{typ} validation, the signature would
verify and many claim checks would pass. The \texttt{typ: "at+jwt"} header is the
mechanism defined by 
\href{https://datatracker.ietf.org/doc/html/rfc9068#section-2.1}{RFC~9068~\S2.1} 
specifically to defeat this class of attack. The \texttt{typ} check MUST be 
performed before any payload claim is inspected.



\textbf{Cognito exception.} AWS Cognito emits \texttt{typ: "JWT"} and
\texttt{token\_use: "access"} instead of \texttt{typ: "at+jwt"}. This is a documented
non-compliance with \href{https://datatracker.ietf.org/doc/html/rfc9068}{RFC~9068}. Canvasflow accommodates this via a per-tenant
\texttt{typ\_strict} configuration flag. When a tenant's configuration carries
\texttt{typ\_strict = false} (set at onboarding for Cognito tenants), the Resource
Server accepts \texttt{typ: "JWT"} and enforces \texttt{token\_use == "access"} as a
substitute check. Tenants using Cognito MUST declare this at onboarding. All other
tenants MUST produce \texttt{typ: "at+jwt"}.

\subsubsection{\texttt{alg} --- Signing Algorithm (RFC~8725~\S3.1)}

The \texttt{alg} header MUST be one of: \texttt{"RS256"}, \texttt{"ES256"},
\texttt{"HS256"}.



The value \texttt{"none"} MUST be rejected unconditionally regardless of any other token
properties (\href{https://datatracker.ietf.org/doc/html/rfc8725#section-3.1}{RFC~8725~\S3.1}). 
The Resource Server MUST validate the \texttt{alg} header against the allowlist 
before performing any verification and MUST NOT use an algorithm
value taken from the token without this check.

\subsubsection{\texttt{kid} --- Key Identifier}

The \texttt{kid} header MUST be present. For JWKS tenants, it selects the public key
from the JWKS endpoint. For HS256 tenants, it selects the secret version from the stored
secret set. A token without \texttt{kid} MUST be rejected.

\pagebreak

\subsection{Required Payload Claims}\label{sec:required-payload-claims}

The following claims MUST be present in the JWT payload of every access token:

\begin{longtable}{@{}lp{1.5cm}p{8cm}@{}}
\toprule
\textbf{Claim} & \textbf{Type} & \textbf{Requirement} \\
\midrule
\endhead
\texttt{iss} & string &
  Tenant issuer URL. MUST exactly match the value registered in \texttt{tenant\_config}.
  Trailing slash is significant: \texttt{https://login.tenant.com} $\neq$
  \texttt{https://login.tenant.com/}. \\[6pt]
\texttt{sub} & string &
  Stable, opaque user identifier assigned by the tenant's IdP. MUST NOT be the user's
  email address (\href{https://datatracker.ietf.org/doc/html/rfc9068\#section-2.2}{RFC~9068~\S2.2}). 
  The \texttt{sub} value MUST remain stable across
  sessions for the same user. \\[6pt]
\texttt{aud} & string or array &
  MUST contain \texttt{"canvasflow-api"}. MAY be an array of strings
  (\href{https://datatracker.ietf.org/doc/html/rfc7519\#section-4.1.3}{RFC~7519~\S4.1.3} 
  permits arrays). Validation passes if \texttt{"canvasflow-api"} is
  \emph{present} in the \texttt{aud} value --- it need not be the only value. \\[6pt]
\texttt{iat} & number &
  Unix timestamp of issuance (seconds since epoch). MUST NOT be more than 60~seconds in
  the future (clock skew tolerance). \\[6pt]
\texttt{exp} & number &
  Unix timestamp of expiry (seconds since epoch). MUST be greater than the current time
  (with $\leq$~60 seconds clock skew tolerance). The difference
  $\texttt{exp} - \texttt{iat}$ MUST NOT exceed 3600~seconds. \\[6pt]
\texttt{jti} & string &
  Unique token identifier. A UUID is RECOMMENDED (\href{https://datatracker.ietf.org/doc/html/rfc4122}{RFC~4122}). 
  Used for server-side log correlation; NOT used for revocation (by design --- 
  see \hyperref[sec:canvasflow-responsibilities]{\S3.3} and 
  \hyperref[sec:out-of-scope]{\S1.4}). \\[6pt]
Entitlement claim & array &
  One of \texttt{cf:entitlements}, \texttt{resources}, or \texttt{entitlements} per the
  precedence rules in \hyperref[sec:claim-name-hierarchy]{\S4.3}. The resolved value MUST
  be a JSON array of strings. An empty array is valid. \\
\bottomrule
\end{longtable}

\begin{quote}
  \textbf{Auth0 \texttt{aud} note.} When the \texttt{openid} scope is requested, Auth0
  includes both \linebreak \texttt{"canvasflow-api"} and the Auth0 \texttt{/userinfo} URL in the
  \texttt{aud} claim as an array. This is fully compliant with 
  \href{https://datatracker.ietf.org/doc/html/rfc7519\#section-4.1.3}{RFC~7519~\S4.1.3}.
  Canvasflow's validation passes as long as \texttt{"canvasflow-api"} is present in the
  array.
\end{quote}

\pagebreak

\subsection{Prohibited}\label{sec:prohibited}

The following MUST NOT appear in any conforming access token:





\begin{longtable}{@{}p{5.5cm}p{9cm}@{}}
\toprule
\textbf{Prohibited} & \textbf{Reason} \\
\midrule
\endhead
\texttt{alg: "none"} &
  Eliminates signature verification entirely; unconditional rejection per 
  \href{https://datatracker.ietf.org/doc/html/rfc8725\#section-3.1}{RFC~8725~\S3.1}. \\[6pt]
\texttt{typ} absent, \texttt{typ: "JWT"} (except Cognito with \texttt{typ\_strict=false}), \texttt{typ: "id\_token+jwt"}, or any other non-compliant \texttt{typ} &
  Enables token confusion attacks (\href{https://datatracker.ietf.org/doc/html/rfc9068\#section-2.1}{RFC~9068~\S2.1}). \\[6pt]
PII in the access token payload (name, email, physical address, phone number) &
  Access tokens traverse multiple systems; PII creates unnecessary exposure risk. \\[6pt]
RSA key shorter than 2048~bits &
  Computationally feasible to break; below the minimum security level for RS256
  (\href{https://datatracker.ietf.org/doc/html/rfc7518\#section-3.3}{RFC~7518~\S3.3}). \\[6pt]
$\texttt{exp} - \texttt{iat}$ greater than 3600~seconds &
  Exceeds the maximum allowed revocation window (\hyperref[sec:revocation]{\S1.4}, Core Design Principle~3). \\[6pt]
Entitlement claim placed in the ID token instead of the access token &
  The ID token is consumed by the SPA and never sent to the API; the claim would be
  silently absent from the Resource Server's perspective. \\[6pt]
Email address as \texttt{sub} &
  Email addresses are mutable; \texttt{sub} must be a stable, opaque identifier
  (\href{https://datatracker.ietf.org/doc/html/rfc9068\#section-2.2}{RFC~9068~\S2.2}). \\
\bottomrule
\end{longtable}

\pagebreak

\subsection{Conforming Token Example}

The following example shows a fully conforming JWT access token using the JWKS (RS256)
signing path and the recommended \texttt{cf:entitlements} claim name.
\linebreak\linebreak
\textbf{Header:}

\begin{lstlisting}[language=json]
{
  "typ": "at+jwt",
  "alg": "RS256",
  "kid": "tenant-key-v3"
}
\end{lstlisting}

\textbf{Payload:}

\begin{lstlisting}[language=json]
{
  "iss": "https://login.tenant.com",
  "sub": "usr_7f3a9c1d-opaque",
  "aud": "canvasflow-api",
  "iat": 1720224000,
  "exp": 1720227600,
  "jti": "a8f3c1d2-9e4b-4f7a-b6c5-1234567890ab",
  "client_id": "canvasflow-pwa",
  "cf:entitlements": ["cf:issue:123", "cf:issue:201"]
}
\end{lstlisting}



Notes on this example:
\begin{itemize}
  \item $\texttt{exp} - \texttt{iat} = 3600$ seconds --- the maximum allowed.
  \item \texttt{cf:entitlements} is used (recommended claim name --- works on all IdPs).
  \item All entitlement values carry the mandatory \texttt{cf:} prefix.
  \item \texttt{sub} is an opaque identifier, not an email address.
  \item \texttt{jti} is a UUID per \href{https://datatracker.ietf.org/doc/html/rfc4122}{RFC~4122}.
\end{itemize}

% The full ordered validation sequence the Resource Server applies to this token is defined in \texttt{spec\_appendix\_validation}.
